\documentclass[11pt]{article}
\usepackage[margin=1.05in]{geometry}
\usepackage{amsmath,amssymb,amsthm,mathtools}
\usepackage{microtype}
\usepackage{enumitem}
\usepackage[colorlinks=true,linkcolor=blue,citecolor=blue,urlcolor=blue]{hyperref}

\newtheorem{theorem}{Theorem}[section]
\newtheorem{lemma}[theorem]{Lemma}
\newtheorem{proposition}[theorem]{Proposition}
\newtheorem{corollary}[theorem]{Corollary}
\newtheorem{remark}[theorem]{Remark}

\newcommand{\one}{\mathbf 1}
\newcommand{\R}{\mathbb R}
\newcommand{\eps}{\varepsilon}
\DeclareMathOperator{\Tr}{Tr}

\title{Partial $C_{13}$ progress: \\
\large path-certificate range, a closed-form linear polynomial, and the marginal gap}
\author{Working note for the odd-cycle project}
\date{May 30, 2026}

\begin{document}
\maketitle

\begin{abstract}
We extend the complement-expansion path-certificate of the $C_9$ and $C_{11}$ proofs to $C_{13}$, establishing
\[
        t(C_{13}, W) \ge p^{13} - p(1-p)^{12}
\]
for every graphon $W$ with edge density $p \ge 519/1000$, by an exact rational Bernstein-subdivision certificate.  As a byproduct we identify a closed-form formula for the linear-part polynomial $P_q^{(m)}(\lambda)$ valid uniformly for all odd $m \ge 5$:
\[
        P_q^{(m)}(\lambda) = \sum_{j=0}^{m-3} a_j^{(m)}(q)\, \lambda^j,
        \qquad
        a_j^{(m)}(q) = (-1)^j m (1-q)^{m-2-j} + m q^{m-2-j} - (m-2-j) q^{m-3-j}.
\]
The full $C_{13}$ case is \emph{not} closed by the hybrid strategy of $C_9$ and $C_{11}$: the spectral-triangle near-bipartite argument generalises to $C_{13}$ but only covers $\eps = p - 1/2 \in (0, \delta_{13}]$ with $\delta_{13} = 0.01559\ldots$, strictly less than $\rho_{13} - 1/2 = 19/1000$.  A residual uncovered window $p \in (1/2 + \delta_{13},\, 519/1000)$ of width $\approx 0.0034$ remains.  This note documents the rigorous path-certificate portion and the closed-form formula, both of which appear new, and locates precisely where the hybrid strategy first fails.
\end{abstract}

\section{Statement and partial result}

The project conjecture is
\begin{equation}\label{eq:conj}
        t(C_{2k+1}, W) \ge p^{2k+1} - p(1-p)^{2k}
\end{equation}
for every graphon $W$ with $p = t(K_2, W)$.  The cases $C_3, C_5, C_7, C_9, C_{11}$ are settled by the main paper.  This note establishes:

\begin{theorem}[Path-certificate range for $C_{13}$]\label{thm:C13-path}
For every graphon $W$ with $p = t(K_2, W) \ge 519/1000$,
\[
        t(C_{13}, W) \ge p^{13} - p(1-p)^{12}.
\]
\end{theorem}

For $p \le 1/2$ the proposed right-hand side is non-positive, so \eqref{eq:conj} is trivial.  Theorem~\ref{thm:C13-path} therefore reduces the open range to $1/2 < p < 519/1000$; the spectral-triangle argument of \S\ref{sec:spectral} covers $p \le 1/2 + \delta_{13} \approx 0.5156$, leaving the narrow residual $p \in (0.5156, 0.519)$ open.

\section{The complement-expansion path certificate for $C_{13}$}\label{sec:path}

We follow the framework of \S 6 (for $C_9$) and \S 7 (for $C_{11}$) in the main paper.  Put $U = 1 - W$ and $q = t(K_2, U) = 1 - p$.  For the path with $j$ edges write $x_j = t(P_j, U)$, and let $c_{13} = t(C_{13}, U)$.  Expanding $t(C_{13}, 1-U)$ by inclusion-exclusion over the thirteen edges of $C_{13}$ and using the deletion bound $c_{13} \le x_{12}$ gives
\begin{equation}\label{eq:Phi13-defn}
        t(C_{13}, 1-U) - \bigl((1-q)^{13} - (1-q) q^{12}\bigr) \ge \Phi_{13},
\end{equation}
where $\Phi_{13}$ is an explicit polynomial in $q, x_2, \ldots, x_{12}$ with $259$ monomials.  The accompanying script \texttt{odd\_cycle\_c13\_checker.py} produces $\Phi_{13}$ by inclusion-exclusion in $<1$ s.  The polynomial is too large to display in print, but the next subsection's moment decomposition is compact and self-contained.

\subsection{Moment decomposition}

Let $T = T_U$.  Write $d = T\one = q\one + g$ with $\langle g, \one \rangle = 0$, and let
\[
        A = P_{\one^\perp} T P_{\one^\perp}
\]
be the compression of $T$ to $\one^\perp$.  Define
\[
        s_j = \langle g, A^j g \rangle, \qquad j \ge 0,
\]
so $s_0 = \|g\|_2^2 \ge 0$.  The path densities $x_j = \langle \one, T^j \one\rangle$ are obtained from the block recurrence
\[
        T(a\one + h) = (aq + \langle g, h\rangle)\one + (ag + Ah), \qquad h \in \one^\perp,
\]
extended to $j = 12$.  Substituting the resulting formulae into $\Phi_{13}$ gives the homogeneous decomposition
\begin{equation}\label{eq:Phi13-decomp}
        \Phi_{13} = L_1 + L_2 + L_3 + L_4 + L_5 + 12\, s_0^6,
\end{equation}
where $L_d$ is homogeneous of degree $d$ in $s_0, s_1, \ldots, s_{10}$.  The pure degree-$6$ term $12 s_0^6 \ge 0$ is immediate.

\begin{remark}[Pattern in the top term]
The high-degree pure term in $\Phi_m$ for odd $m$ is $(m-1)\, s_0^{(m-1)/2}$: for $m = 11$ this is $10 s_0^5$ (\S 7.2 of the main paper), and for $m = 13$ it is $12 s_0^6$.  The pattern is verified by the checker for $m = 13$ and is consistent with the closed-form leading coefficient $a_{m-3}^{(m)}(q) \equiv m - 1$ established in Proposition~\ref{prop:Pq-closed} below.
\end{remark}

\subsection{The linear part and the closed-form $P_q^{(m)}$}

The degree-$1$ part is
\[
        L_1 = \int_{-1}^{1} P_q^{(13)}(\lambda)\, d\mu(\lambda),
\]
where $\mu$ is the spectral measure of $A$ with respect to $g$ and $P_q^{(13)}$ is a polynomial of degree $m - 3 = 10$ in $\lambda$.

A central observation, verified symbolically by the accompanying script, is that the polynomial $P_q^{(m)}(\lambda)$ produced in this way admits a uniform closed-form expression for every odd $m \ge 5$.

\begin{proposition}[Closed form for $P_q^{(m)}$]\label{prop:Pq-closed}
For every odd integer $m \ge 5$, the linear-part polynomial in the moment decomposition of $\Phi_m$ is
\begin{equation}\label{eq:Pq-closed}
        P_q^{(m)}(\lambda) = \sum_{j=0}^{m-3} a_j^{(m)}(q)\, \lambda^j,
\end{equation}
with coefficients
\begin{equation}\label{eq:a-coeffs}
        a_j^{(m)}(q)
        = (-1)^j\, m\, (1-q)^{m-2-j}
        + m\, q^{m-2-j}
        - (m-2-j)\, q^{m-3-j}
\end{equation}
for $0 \le j \le m - 3$.  In particular the leading coefficient $a_{m-3}^{(m)}(q) = m - 1$.
\end{proposition}

The proposition is verified by the checker at $m = 11$ (matching the explicit $P_q^{(11)}$ of \S 7.2 in the main paper) and at $m = 13$ (matching the inclusion-exclusion-derived $P_q^{(13)}$).  Specialising \eqref{eq:Pq-closed}--\eqref{eq:a-coeffs} at $m = 13$:
\begin{align*}
P_q^{(13)}(\lambda)
={}& 12\lambda^{10} + (24q - 13)\lambda^9 + (36q^2 - 39q + 13)\lambda^8 \\
&{}+ (48q^3 - 78q^2 + 52q - 13)\lambda^7 + (60q^4 - 130q^3 + 130q^2 - 65q + 13)\lambda^6 \\
&{}+ (72q^5 - 195q^4 + 260q^3 - 195q^2 + 78q - 13)\lambda^5 \\
&{}+ (84q^6 - 273q^5 + 455q^4 - 455q^3 + 273q^2 - 91q + 13)\lambda^4 \\
&{}+ (96q^7 - 364q^6 + 728q^5 - 910q^4 + 728q^3 - 364q^2 + 104q - 13)\lambda^3 \\
&{}+ (108q^8 - 468q^7 + 1092q^6 - 1638q^5 + 1638q^4 - 1092q^3 + 468q^2 - 117q + 13)\lambda^2 \\
&{}+ (120q^9 - 585q^8 + 1560q^7 - 2730q^6 + 3276q^5 - 2730q^4 + 1560q^3 - 585q^2 + 130q - 13)\lambda \\
&{}+ 132q^{10} - 715q^9 + 2145q^8 - 4290q^7 + 6006q^6 - 6006q^5 + 4290q^4 - 2145q^3 + 715q^2 - 143q + 13.
\end{align*}

\begin{remark}[Specialisations to $m = 9, 11$]
At $m = 9$ formula~\eqref{eq:a-coeffs} reproduces the $P_q^{(9)}$ from \S 6.1 of the main paper (cf.\ eq.~(6.7) there); at $m = 11$ it reproduces the $P_q^{(11)}$ from \S 7.2 (cf.\ eq.~(7.5) there).  The checker cross-verifies both specialisations.
\end{remark}

\subsection{Nonlinear kernels}

If $s_0 > 0$, let $\nu = \mu/s_0$ be the normalised spectral measure on $[-1, 1]$ and set $m_j = \int_{-1}^1 \lambda^j\, d\nu(\lambda)$.  Then for $d = 1, 2, 3, 4, 5$,
\[
        L_d = s_0^d\, F_d(q;\, m_1, \ldots, m_{10}),
\]
and each $F_d$ is the $d$-fold $\nu$-average of an explicit symmetric kernel $K_{d, q}$ on $[-1, 1]^d$:
\[
        F_d = \int_{[-1, 1]^d} K_{d, q}(\lambda_1, \ldots, \lambda_d)\, d\nu(\lambda_1) \cdots d\nu(\lambda_d).
\]
The kernel is obtained by replacing each monomial $m_{a_1} \cdots m_{a_t}$ in $F_d$ by the symmetrised product
\[
        \frac{1}{d!} \sum_{\sigma \in S_d} \lambda_{\sigma(1)}^{a_1} \cdots \lambda_{\sigma(t)}^{a_t}
\]
with the remaining $d - t$ exponents equal to $0$, as in \S 7.2 of the main paper.

\subsection{Bernstein-certified positivity}

The accompanying script certifies, by exact rational Bernstein subdivision, the following positivity statements:
\begin{align}
P_q^{(13)}(\lambda) > 0 & \quad \text{on } q \in [0, 481/1000],\ \lambda \in [-1, 1], \label{eq:Pq-positive-C13} \\
K_{2, q}(\lambda_1, \lambda_2) > 0 & \quad \text{on } q \in [0, 481/1000],\ \lambda_i \in [-1, 1], \label{eq:K2-positive-C13} \\
K_{d, q}(\lambda_1, \ldots, \lambda_d) > 0 & \quad \text{on } q \in [0, 1/2],\ \lambda_i \in [-1, 1], \quad d = 3, 4, 5. \label{eq:Kd-positive-C13}
\end{align}
The subdivision certificates are small: $6$ boxes for $P_q^{(13)}$, $6$ for $K_2$, $8$ for $K_3$, $16$ for $K_4$, and $42$ for $K_5$; the full set runs in $\approx 60$ s in sympy.

\begin{remark}[Tightness of $q \le 481/1000$]\label{rmk:tight}
At $q = 482/1000$ the linear polynomial $P_q^{(13)}(\lambda)$ attains negative values near $\lambda \approx 0.355$ (a numerical sweep confirms a local minimum below zero), so the threshold $q = 481/1000$ for \eqref{eq:Pq-positive-C13} is tight up to the rational denominator $1000$.  The nonlinear kernel $K_2$ likewise dips slightly negative near $q = 1/2$ for $C_{13}$, requiring the same $q \le 481/1000$ cap.  This is in contrast to $C_{11}$, where $K_2$ is positive on the full $q \in [0, 1/2]$.
\end{remark}

\subsection{Conclusion of the path-certificate range}

Combining the decomposition \eqref{eq:Phi13-decomp} with the positivity statements \eqref{eq:Pq-positive-C13}, \eqref{eq:K2-positive-C13}, \eqref{eq:Kd-positive-C13} gives
\[
        \Phi_{13} = L_1 + L_2 + L_3 + L_4 + L_5 + 12\, s_0^6 \ge 0
        \qquad \text{on } 0 \le q \le 481/1000.
\]
By \eqref{eq:Phi13-defn} this yields Theorem~\ref{thm:C13-path}, with $\rho_{13} = 519/1000$.

\section{The spectral-triangle near-bipartite window}\label{sec:spectral}

The spectral closure of \S 7.1 in the main paper (originally $C_{11}$ on $1/2 < p \le 103/200$) generalises verbatim to every odd $m \ge 5$.  We sketch the $C_{13}$ specialisation and the size of the resulting window.

Let $T = T_W$ be the self-adjoint Hilbert--Schmidt operator with kernel $W$, with eigenvalues $\lambda_0 \ge 0, \lambda_1, \lambda_2, \ldots$ satisfying $\lambda_0 \ge p$, $\sum_i \lambda_i^2 \le p$, and $t(C_m, W) = \sum_i \lambda_i^m$ for $m \ge 3$.

\begin{lemma}[Spectral closure criterion for $C_{13}$]\label{lem:spectral-C13}
Let $1/2 < p \le 1/2 + \delta_{13}$ and $q = 1 - p$.  Assume that $t(K_3, W) \ge \Theta(p)$, where $\Theta(p) = \tfrac32 c(1-c)^2$ is the Razborov--Reiher sharp triangle bound with $p = \tfrac12 + c - \tfrac32 c^2$.  Then
\[
        t(C_{13}, W) \ge p^{13} - pq^{12}.
\]
\end{lemma}

\begin{proof}[Proof sketch]
The structure is identical to \S 7.1 of the main paper, with the cycle length, exponent, and cut-off relaxed to their $C_{13}$ analogues.  Set $\ell = \lambda_0 \ge p$, $S = p - \ell^2$, $N_{13} = \sum_{i \ge 1,\, \lambda_i < 0} |\lambda_i|^{13}$, and
\[
        \alpha = \alpha_\ell := (\ell^{13} - p^{13} + p q^{12})^{1/13}.
\]
The reduction $G'(\ell) > 0$ on $\Delta = p - \ell^2 - \alpha^2 \ge 0$ goes through because differentiating $\alpha^{13} = \ell^{13} - p^{13} + p q^{12}$ gives $\alpha' = \ell^{12}/\alpha^{12}$ and (dropping the non-negative $\sqrt{\Delta}$ term)
\[
        G'(\ell) > -3\ell^2 + 3 \alpha^2 \alpha' = 3 \ell^2 \Bigl(\bigl(\tfrac{\ell}{\alpha}\bigr)^{10} - 1\Bigr) > 0
\]
whenever $\alpha < \ell$, which holds in our range.  At $\ell = p$ the inequality \eqref{eq:scalar-C13} below is the binding scalar comparison.
\end{proof}

The scalar comparison at $\ell = p$ is
\begin{equation}\label{eq:scalar-C13}
        \Theta(p) \ge p^3 - \alpha_0^3 + \Delta_0^{3/2},
        \qquad \alpha_0 = (p q^{12})^{1/13},\quad \Delta_0 = pq - \alpha_0^2.
\end{equation}

\begin{proposition}[Maximal spectral window]\label{prop:delta13}
The maximal $\delta > 0$ such that \eqref{eq:scalar-C13} holds on $\eps = p - 1/2 \in (0, \delta]$ is
\[
        \delta_{13} = 0.015592935063\ldots,
\]
computed by bisection in $50$-digit arithmetic.  In particular $\delta_{13} < 19/1000 = \rho_{13} - 1/2$.
\end{proposition}

The bisection is documented in the synthesis report \texttt{general\_odd\_cycle\_progress.md} accompanying this directory.  An asymptotic Taylor analysis there shows $m^2 \delta_m \to 9/4$ as $m \to \infty$; for $m = 13$ this gives the leading-order estimate $9 \cdot 13 / (4 \cdot 11^3) \approx 0.0266$, with the actual $\delta_{13} = 0.01559$ being smaller due to $\eps^2$ and higher corrections.

\section{The marginal gap and what remains}\label{sec:remains}

\subsection{The uncovered residual interval}

Combining the trivial case, Lemma~\ref{lem:spectral-C13}/Proposition~\ref{prop:delta13}, and Theorem~\ref{thm:C13-path}, the $C_{13}$ conjecture is established for
\[
        p \le 1/2 \quad \text{(trivially)}, \qquad 1/2 < p \le 1/2 + \delta_{13} \approx 0.5156 \quad \text{(spectral)}, \qquad p \ge 519/1000 \quad \text{(path-certificate)}.
\]
The residual uncovered interval is
\[
        p \in \bigl(1/2 + \delta_{13},\ 519/1000\bigr) = (0.51559\ldots,\ 0.519),
\]
of width $\approx 0.00341$.  This is the first cycle length at which the hybrid strategy of \S 6 and \S 7 fails to close the conjecture.

\subsection{Plausible routes to closing the gap}

\begin{enumerate}[leftmargin=2em]
\item \emph{Tighten the path-certificate by a $\sim 25\%$ improvement.}  Either of the following would close the residual:
\begin{itemize}
\item A sharper Bernstein certificate for $K_{2, q}$ that pushes the $q$-ceiling from $481/1000$ down to $\approx 1/2 - \delta_{13}$.
\item A refined positivity certificate for $P_q^{(13)}(\lambda)$ that extends $\rho_{13}$ further toward $1/2$.
\end{itemize}
Both seem feasible: the binding constraint at $q = 482/1000$ is local (near $\lambda \approx 0.355$) and the negative dip is shallow.  A non-Bernstein certificate (e.g.\ an SDP, or a tighter sum-of-squares) may close it.

\item \emph{Refine the spectral closure to $O(\eps^2)$.}  The $\eps^2$ corrections to \eqref{eq:scalar-C13} were dropped in the leading-order analysis behind Proposition~\ref{prop:delta13}; including them adds a positive correction of order $1/m$ to the slack, potentially raising $\delta_{13}$ by a few tens of percent.

\item \emph{Stronger near-bipartite input.}  A sharp $t(C_5, W) \ge t(C_5, T_{2;c})$ Tur\'an-type bound on $p \in [1/2, 2/3]$ (currently open) would replace the triangle-density input in Lemma~\ref{lem:spectral-C13} by a substantially stronger one, giving $\delta_m \ge 0.122$ \emph{uniformly in $m$} and closing every odd cycle simultaneously, provided the path-certificate cut-off satisfies $\rho_m \le 1/2 + 0.122$.
\end{enumerate}

\subsection{Asymptotic obstruction}

Independent of $C_{13}$ specifically, Taylor analysis of \eqref{eq:scalar-C13} shows the present spectral-triangle window has a hard ceiling
\[
        \delta_m \le \frac{9 m}{4 (m-2)^3} \sim \frac{9}{4 m^2}
        \qquad (m \to \infty).
\]
This is intrinsic to using triangle density together with the $L^2$-to-$L^3$ comparison underlying \eqref{eq:scalar-C13}: no manipulation within these inputs can change the $1/m^2$ scaling.  Meanwhile the path-certificate cut-offs $\rho_m - 1/2$ have grown rapidly ($3/2000$ at $C_9$, $3/200$ at $C_{11}$, $19/1000$ at $C_{13}$), and numerical experiments (in the synthesis report) suggest they may stabilise around a constant $\approx 0.02$ for large $m$ rather than shrinking.  In particular, the two windows divergently miss each other from $C_{13}$ onward unless one of the strengthenings in \S\ref{sec:remains}.2 is achieved.

\section*{Accompanying artefacts}

The accompanying script \texttt{odd\_cycle\_c13\_checker.py} in the same directory verifies, in exact rational arithmetic:
\begin{itemize}[leftmargin=2em]
\item the degree-$0$ and degree-$6$ structural parts of $\Phi_{13}$;
\item the closed-form Proposition~\ref{prop:Pq-closed} at $m = 11$ (cross-check) and $m = 13$;
\item the positivity statements \eqref{eq:Pq-positive-C13}--\eqref{eq:Kd-positive-C13} by exact rational Bernstein subdivision.
\end{itemize}
The script completes in $\approx 60$ s.

The spectral-window value $\delta_{13} = 0.015592935063\ldots$ in Proposition~\ref{prop:delta13} is not verified by this checker; it is documented in the synthesis report \texttt{general\_odd\_cycle\_progress.md}, produced by a separate $50$-digit-precision bisection in mpmath.

\section{Addendum: closure attempt on the residual gap (negative)}\label{sec:closure-attempt}

A dedicated attempt (May 31, 2026) to close the residual band
$q\in(0.481,\,0.48441)$, i.e.\ $p\in(0.51559,0.519)$, i.e.\ $\eps=p-\tfrac12\in(0.01559,0.019)$,
by four independent rigorous strategies did \textbf{not} succeed.  Each natural single-method
route was shown to fall provably short, and the obstruction is the same coupling bottleneck as
the general conjecture (the ``aggregate-not-modewise'' spectral bottleneck of the main paper),
here localised to a single frontier eigenvalue.  The findings, all independently re-verified, are recorded because
several are sharp and reusable.

\paragraph{The spectral closure has a hard ceiling $\delta_{13}=0.01559<0.019$.}
Even after adjoining the support bound (next paragraph), the Sidorenko bound
$t(C_4,W)\ge p^4$, and the sharp triangle bound, the maximal window for the spectral-moment method
is \emph{exactly} $\delta_{13}=0.0155929\ldots$ (the scalar quantity $\Theta(p)-(p^3-\alpha_0^3+
\Delta_0^{3/2})$ is $+1.9\times10^{-7}$ at $\eps=0.01559$ but $-2.6\times10^{-4}$ at $\eps=0.019$).
At the binding configuration the negative atom sits at $-\alpha_0$ with $\alpha_0=(pq^{12})^{1/13}
\in(0.4838,0.4867)$ strictly inside $(-\tfrac12,\tfrac12)$, so the support bound is \emph{slack} and
does not help.  A clean fact survives, sharp but inactive here:
\[
        \lambda_{\min}(T_W)\ge-\tfrac12\quad\text{for every graphon }W,
\]
via $T_W=\tfrac12(J+T_S)$ with $S=2W-1$, $\|T_S\|\le1$ by Schur's test (equality at balanced
complete bipartite).

\paragraph{The path certificate provably fails inside the gap.}
The defect polynomial $\Phi_{13}$ of \S\ref{sec:path} is \emph{not} nonnegative on the
realisable region for $q$ in the gap.  An exact rational two-block witness with \emph{unequal}
part measures $(39/100,\,61/100)$ and complement-block values
$U=\bigl(\begin{smallmatrix}43/50&1/10\\1/10&41/50\end{smallmatrix}\bigr)$ gives
$q=120877/250000=0.483508$ (inside the gap) and
\[
        \Phi_{13}=-2.0758\times10^{-7}<0,
\]
while the true defect is $t(C_{13},W)-g_{13}(p)=+1.105\times10^{-4}>0$.  The entire positive margin
lives in the discarded deletion term $x_{12}-c_{13}=t(P_{12},U)-t(C_{13},U)=+1.107\times10^{-4}$.
Hence the path-certificate range cannot be pushed below $\rho_{13}=519/1000$ \emph{without retaining
the deletion gap}.  (This corrects an earlier, over-optimistic diagnostic that had restricted
attention to equal-measure blocks, where $\Phi_{13}\ge0$ does persist further.)

\paragraph{A useful new spectral lemma.}
Since $\Tr(A^2)\le pq$ (from $\int U^2\le\int U$), and two eigenvalues exceeding $q$ would force
$\Tr(A^2)>2q^2>pq$ whenever $q>1/3$, the compression $A=P_{\one^\perp}T_UP_{\one^\perp}$ has
\emph{at most one} eigenvalue $\alpha_1>q$ in the nontrivial range.  Consequently
\[
        \Tr(A^{13})-pq^{12}\le \alpha_1^2\bigl(\alpha_1^{11}-q^{11}\bigr),
        \qquad \alpha_1\in(q,\tfrac12],
\]
localising the entire obstruction to one frontier eigenvalue.  However, the modewise
hub-coupling surplus cannot pay for this excess: the coupling obeys $b_1^2\le p(\tfrac12-\alpha_1)$
(so the surplus $\to0$ exactly as $\alpha_1\to\tfrac12$, where the excess is largest), and an exact
rational witness exhibits a frontier mode $\alpha_1=0.4885>q$ that is \emph{fully decoupled}
($b_1=\langle g,\phi_1\rangle=0$ by a swap symmetry) while the true defect stays positive.  The required
surplus must therefore come from the aggregate, not from the frontier mode --- the
``aggregate-not-modewise'' bottleneck.

\paragraph{Status and the route forward.}
The residual $C_{13}$ band remains open.  All evidence is that the conjecture holds there (every
search yields defect $\gtrsim+10^{-5}$), but no marginal-moment certificate is strong enough,
because the minimising spectral configuration of every relaxation is an \emph{unrealisable} spectrum
that the $g$--$A$ coupling forbids.  The single most promising route is a \emph{deletion-retaining}
certificate: lower-bound $x_{12}-c_{13}$ in the moment variables $s_0,\dots,s_{10}$ and add it back to
$\Phi_{13}$, then seek a sum-of-squares certificate for $\Phi_{13}+L(s_0,\dots,s_{10})\ge0$ on the band,
using the one-frontier-mode reduction above.  The unconditional fallback is the (open) sharp
$C_5$ Tur\'an-type lower bound on $[\tfrac12,\tfrac23]$, which would overshoot the band entirely.

\begin{thebibliography}{9}

\bibitem{RazborovTriangles}
A.~A. Razborov.
\newblock On the minimal density of triangles in graphs.
\newblock \emph{Combinatorics, Probability and Computing} 17(4):603--618, 2008.

\bibitem{ReiherCliqueDensity}
C.~Reiher.
\newblock The clique density theorem.
\newblock \emph{Annals of Mathematics} 184(3):683--707, 2016.

\end{thebibliography}

\end{document}
